\documentclass[12pt]{article}

\usepackage{epsfig}
\usepackage[spanish,activeacute]{babel}
\usepackage[latin1]{inputenc}
\usepackage[reqno]{amsmath}
\usepackage{graphicx}
\usepackage{amsfonts}
\usepackage{mathrsfs}
\usepackage{amssymb}
\usepackage{latexsym}

\usepackage[usenames,dvipsnames]{pstricks}
%\usepackage{epsfig}
\usepackage{pst-grad} % For gradients
\usepackage{pst-plot} % For axes

\topmargin -2.5 cm
\oddsidemargin -1.5 cm
\textheight 26.5 cm
\textwidth 18.5cm

\newcommand{\abre}{\textquestiondown}

\begin{document}
\pagestyle{empty}
\sffamily
%\twocolumn[

%\begin{minipage}{5 cm}{\begin{center}
%\includegraphics[scale=0.7]{escudo}\end{center}}\end{minipage}
%\begin{minipage}{11 cm}{
%\begin{flushleft}\vspace*{0.2cm} \textbf{Universidad de Antioquia}\\[2mm] 
%Facultad de Ciencias Exactas y Naturales\\[2mm] 
%Instituto de F�sica\\[2mm] 
F�sica de campos. TALLER 9 - \textbf{Ley de Ampere} 
%\end{flushleft}
%}
%\end{minipage}


\hrulefill
%\footnote{Algunas de las figuras son tomadas del texto: \textit{Physics For Scientists And Engineers 6E By Serway And Jewett.pdf}} 

\vspace{0.5cm}
%]











{\color{red}LEY DE AMPERE}
\vspace{0.6 cm}

\begin{enumerate} 

\item Leidy Usando la ley de Ampere, calcule el campo magn�tico en el interior de solenoide muy largo de longitud $l$, de $N$ vueltas y de corriente $I$.
\begin{displaymath}
B=\mu_0 \dfrac{N}{l}I = \mu_0 n I \hspace*{0.5 cm};\hspace*{0.5 cm} \dfrac{N}{l}=n
\end{displaymath}

\item Usando la ley de Ampere, calcule el campo magn�tico dentro de un toroide e radio $r$.
\begin{displaymath}
B=\mu_0 \dfrac{NI}{2\pi r}
\end{displaymath}

\item A lo largo de un cilindro de radio $R$ fluye una corriente $I$ distribuida de modo uniforme. Calcule el campo magn�tico $\vec{B}$ en el interior y en el exterior del cable.
\begin{displaymath}
B=\dfrac{\mu_0}{2\pi}\dfrac{r}{R^2} \hspace*{1 cm} 0\ll r \ll R
\end{displaymath}
\begin{displaymath}
B=\dfrac{\mu_0}{2\pi} \dfrac{1}{r} \hspace*{1 cm} r\gg R
\end{displaymath}

\item Calcule el campo magn�tico de un cilindro de radio $R$ por el que fluye una corriente cuya densidad es $J=br$, donde $b$ es una constante y $r$ es la distancia radial al centro del cilindro.
\begin{displaymath}
B=\mu_0 b \dfrac{r^2}{3} \hspace*{1 cm} 0\ll r\ll R
\end{displaymath}
\begin{displaymath}
B=\mu_0 b \dfrac{R^3}{3}\dfrac{1}{r} \hspace*{1 cm} r\gg R
\end{displaymath}

\item Se tienen un cilindro de radio $R_{1}$ conc�ntrico con un cascar�n cil�ndrico de radio interno $R_{2}$ y radio externo $R_{3}$ tal como se muestra en la figura. El cilindro interior tiene una densidad de corriente $\vec{J}_{1}=b r \hat{k}$, donde $b$ es una constante y $r$ es la distancia al centro. El cascar�n cil�ndrico tiene una densidad de corriente constante $\vec{J}_{2}=-J_{2}\hat{k}$.\\

\begin{minipage}{7 cm}{
Hallar el campo magn�tico para:
\begin{enumerate}
\item $r \le R_{1}$
\item $R_{1} \le r \le R_{2}$
\item $R_{2} \le r \le R_{3}$
\item $r \geq R_{3}$
\end{enumerate}
}
\end{minipage}
\begin{minipage}{10 cm}{
\includegraphics[scale=0.6]{cilindro}
}
\end{minipage}

\vspace*{1cm}













\end{enumerate}
\end{document}
